% Options for packages loaded elsewhere
% Options for packages loaded elsewhere
\PassOptionsToPackage{unicode}{hyperref}
\PassOptionsToPackage{hyphens}{url}
\PassOptionsToPackage{dvipsnames,svgnames,x11names}{xcolor}
%
\documentclass[
  11pt,
  a4paper,
  DIV=11,
  numbers=noendperiod]{scrartcl}
\usepackage{xcolor}
\usepackage[lmargin=2cm,rmargin=2cm,tmargin=1.5cm,bmargin=2cm]{geometry}
\usepackage{amsmath,amssymb}
\setcounter{secnumdepth}{-\maxdimen} % remove section numbering
\usepackage{iftex}
\ifPDFTeX
  \usepackage[T1]{fontenc}
  \usepackage[utf8]{inputenc}
  \usepackage{textcomp} % provide euro and other symbols
\else % if luatex or xetex
  \usepackage{unicode-math} % this also loads fontspec
  \defaultfontfeatures{Scale=MatchLowercase}
  \defaultfontfeatures[\rmfamily]{Ligatures=TeX,Scale=1}
\fi
\usepackage{lmodern}
\ifPDFTeX\else
  % xetex/luatex font selection
\fi
% Use upquote if available, for straight quotes in verbatim environments
\IfFileExists{upquote.sty}{\usepackage{upquote}}{}
\IfFileExists{microtype.sty}{% use microtype if available
  \usepackage[]{microtype}
  \UseMicrotypeSet[protrusion]{basicmath} % disable protrusion for tt fonts
}{}
\makeatletter
\@ifundefined{KOMAClassName}{% if non-KOMA class
  \IfFileExists{parskip.sty}{%
    \usepackage{parskip}
  }{% else
    \setlength{\parindent}{0pt}
    \setlength{\parskip}{6pt plus 2pt minus 1pt}}
}{% if KOMA class
  \KOMAoptions{parskip=half}}
\makeatother
% Make \paragraph and \subparagraph free-standing
\makeatletter
\ifx\paragraph\undefined\else
  \let\oldparagraph\paragraph
  \renewcommand{\paragraph}{
    \@ifstar
      \xxxParagraphStar
      \xxxParagraphNoStar
  }
  \newcommand{\xxxParagraphStar}[1]{\oldparagraph*{#1}\mbox{}}
  \newcommand{\xxxParagraphNoStar}[1]{\oldparagraph{#1}\mbox{}}
\fi
\ifx\subparagraph\undefined\else
  \let\oldsubparagraph\subparagraph
  \renewcommand{\subparagraph}{
    \@ifstar
      \xxxSubParagraphStar
      \xxxSubParagraphNoStar
  }
  \newcommand{\xxxSubParagraphStar}[1]{\oldsubparagraph*{#1}\mbox{}}
  \newcommand{\xxxSubParagraphNoStar}[1]{\oldsubparagraph{#1}\mbox{}}
\fi
\makeatother

\usepackage{color}
\usepackage{fancyvrb}
\newcommand{\VerbBar}{|}
\newcommand{\VERB}{\Verb[commandchars=\\\{\}]}
\DefineVerbatimEnvironment{Highlighting}{Verbatim}{commandchars=\\\{\}}
% Add ',fontsize=\small' for more characters per line
\usepackage{framed}
\definecolor{shadecolor}{RGB}{241,243,245}
\newenvironment{Shaded}{\begin{snugshade}}{\end{snugshade}}
\newcommand{\AlertTok}[1]{\textcolor[rgb]{0.68,0.00,0.00}{#1}}
\newcommand{\AnnotationTok}[1]{\textcolor[rgb]{0.37,0.37,0.37}{#1}}
\newcommand{\AttributeTok}[1]{\textcolor[rgb]{0.40,0.45,0.13}{#1}}
\newcommand{\BaseNTok}[1]{\textcolor[rgb]{0.68,0.00,0.00}{#1}}
\newcommand{\BuiltInTok}[1]{\textcolor[rgb]{0.00,0.23,0.31}{#1}}
\newcommand{\CharTok}[1]{\textcolor[rgb]{0.13,0.47,0.30}{#1}}
\newcommand{\CommentTok}[1]{\textcolor[rgb]{0.37,0.37,0.37}{#1}}
\newcommand{\CommentVarTok}[1]{\textcolor[rgb]{0.37,0.37,0.37}{\textit{#1}}}
\newcommand{\ConstantTok}[1]{\textcolor[rgb]{0.56,0.35,0.01}{#1}}
\newcommand{\ControlFlowTok}[1]{\textcolor[rgb]{0.00,0.23,0.31}{\textbf{#1}}}
\newcommand{\DataTypeTok}[1]{\textcolor[rgb]{0.68,0.00,0.00}{#1}}
\newcommand{\DecValTok}[1]{\textcolor[rgb]{0.68,0.00,0.00}{#1}}
\newcommand{\DocumentationTok}[1]{\textcolor[rgb]{0.37,0.37,0.37}{\textit{#1}}}
\newcommand{\ErrorTok}[1]{\textcolor[rgb]{0.68,0.00,0.00}{#1}}
\newcommand{\ExtensionTok}[1]{\textcolor[rgb]{0.00,0.23,0.31}{#1}}
\newcommand{\FloatTok}[1]{\textcolor[rgb]{0.68,0.00,0.00}{#1}}
\newcommand{\FunctionTok}[1]{\textcolor[rgb]{0.28,0.35,0.67}{#1}}
\newcommand{\ImportTok}[1]{\textcolor[rgb]{0.00,0.46,0.62}{#1}}
\newcommand{\InformationTok}[1]{\textcolor[rgb]{0.37,0.37,0.37}{#1}}
\newcommand{\KeywordTok}[1]{\textcolor[rgb]{0.00,0.23,0.31}{\textbf{#1}}}
\newcommand{\NormalTok}[1]{\textcolor[rgb]{0.00,0.23,0.31}{#1}}
\newcommand{\OperatorTok}[1]{\textcolor[rgb]{0.37,0.37,0.37}{#1}}
\newcommand{\OtherTok}[1]{\textcolor[rgb]{0.00,0.23,0.31}{#1}}
\newcommand{\PreprocessorTok}[1]{\textcolor[rgb]{0.68,0.00,0.00}{#1}}
\newcommand{\RegionMarkerTok}[1]{\textcolor[rgb]{0.00,0.23,0.31}{#1}}
\newcommand{\SpecialCharTok}[1]{\textcolor[rgb]{0.37,0.37,0.37}{#1}}
\newcommand{\SpecialStringTok}[1]{\textcolor[rgb]{0.13,0.47,0.30}{#1}}
\newcommand{\StringTok}[1]{\textcolor[rgb]{0.13,0.47,0.30}{#1}}
\newcommand{\VariableTok}[1]{\textcolor[rgb]{0.07,0.07,0.07}{#1}}
\newcommand{\VerbatimStringTok}[1]{\textcolor[rgb]{0.13,0.47,0.30}{#1}}
\newcommand{\WarningTok}[1]{\textcolor[rgb]{0.37,0.37,0.37}{\textit{#1}}}

\usepackage{longtable,booktabs,array}
\usepackage{calc} % for calculating minipage widths
% Correct order of tables after \paragraph or \subparagraph
\usepackage{etoolbox}
\makeatletter
\patchcmd\longtable{\par}{\if@noskipsec\mbox{}\fi\par}{}{}
\makeatother
% Allow footnotes in longtable head/foot
\IfFileExists{footnotehyper.sty}{\usepackage{footnotehyper}}{\usepackage{footnote}}
\makesavenoteenv{longtable}
\usepackage{graphicx}
\makeatletter
\newsavebox\pandoc@box
\newcommand*\pandocbounded[1]{% scales image to fit in text height/width
  \sbox\pandoc@box{#1}%
  \Gscale@div\@tempa{\textheight}{\dimexpr\ht\pandoc@box+\dp\pandoc@box\relax}%
  \Gscale@div\@tempb{\linewidth}{\wd\pandoc@box}%
  \ifdim\@tempb\p@<\@tempa\p@\let\@tempa\@tempb\fi% select the smaller of both
  \ifdim\@tempa\p@<\p@\scalebox{\@tempa}{\usebox\pandoc@box}%
  \else\usebox{\pandoc@box}%
  \fi%
}
% Set default figure placement to htbp
\def\fps@figure{htbp}
\makeatother





\setlength{\emergencystretch}{3em} % prevent overfull lines

\providecommand{\tightlist}{%
  \setlength{\itemsep}{0pt}\setlength{\parskip}{0pt}}



 


\KOMAoption{captions}{tableheading}
\makeatletter
\@ifpackageloaded{caption}{}{\usepackage{caption}}
\AtBeginDocument{%
\ifdefined\contentsname
  \renewcommand*\contentsname{Table of contents}
\else
  \newcommand\contentsname{Table of contents}
\fi
\ifdefined\listfigurename
  \renewcommand*\listfigurename{List of Figures}
\else
  \newcommand\listfigurename{List of Figures}
\fi
\ifdefined\listtablename
  \renewcommand*\listtablename{List of Tables}
\else
  \newcommand\listtablename{List of Tables}
\fi
\ifdefined\figurename
  \renewcommand*\figurename{Figure}
\else
  \newcommand\figurename{Figure}
\fi
\ifdefined\tablename
  \renewcommand*\tablename{Table}
\else
  \newcommand\tablename{Table}
\fi
}
\@ifpackageloaded{float}{}{\usepackage{float}}
\floatstyle{ruled}
\@ifundefined{c@chapter}{\newfloat{codelisting}{h}{lop}}{\newfloat{codelisting}{h}{lop}[chapter]}
\floatname{codelisting}{Listing}
\newcommand*\listoflistings{\listof{codelisting}{List of Listings}}
\makeatother
\makeatletter
\makeatother
\makeatletter
\@ifpackageloaded{caption}{}{\usepackage{caption}}
\@ifpackageloaded{subcaption}{}{\usepackage{subcaption}}
\makeatother
\usepackage{bookmark}
\IfFileExists{xurl.sty}{\usepackage{xurl}}{} % add URL line breaks if available
\urlstyle{same}
\hypersetup{
  pdftitle={EMU Hacettepe EMU660 - Analitik Yöntemlerle Karar Verme 2025 -- 2026 Bahar Dönemi Ara Sınav},
  pdfauthor={Ece Özcan Karatop},
  colorlinks=true,
  linkcolor={blue},
  filecolor={Maroon},
  citecolor={Blue},
  urlcolor={Blue},
  pdfcreator={LaTeX via pandoc}}


\title{EMU Hacettepe\\
EMU660 - Analitik Yöntemlerle Karar Verme\\
2025 -- 2026 Bahar Dönemi\\
Ara Sınav}
\author{Ece Özcan Karatop}
\date{Sınav Tarihi: April 15, 2026}
\begin{document}
\maketitle


\section{myEMU Veritabanı}\label{myemu-veritabanux131}

EMU Hacettepe'nin bölüm başkanı Prof.~Dr.~Murat Caner Testik, mezun
verilerinden yararlanarak bölümün kariyer hizmetlerini geliştirmek
istemektedir. Caner hoca, analitik yöntemlerle karar verme alanındaki
yetkinliğinize güvendiği için yardımınıza başvurmaktadır. Bu doğrultuda,
EMU Hacettepe'nin \texttt{myEMU\ Öğrenci\ ve\ Mezun\ Takip\ Sisteminden}
sağlanan \texttt{myemu\_data} isimli veri setini sizinle paylaşmaktadır.

\begin{itemize}
\item
  \texttt{myemu\_data}, bölüm lisans programından mezun olan ve iş
  deneyimlerine sahip olan mezunlara ait bilgileri içeren bir veri
  setidir.
\item
  Veri seti, mezunların, mezuniyet not ortalamalarına (CGPA), geçmiş iş
  tecrübelerine ve öğrenciyken katıldıkları Erasmus programına dair
  detayları barındırır. Erasmus katılım sütunundaki 1 değeri, ilgili
  kişinin Erasmus programına katıldığını belirtir.
\item
  Bir mezunun, birden fazla iş tecrübesi olabileceği için veri setinde
  birden fazla kaydı (satırı) olabilir. Aynı kişiye ait her bir satırda
  farklı iş tecrübesi olsa da, not ortalaması ve erasmus gibi bilgiler
  aynı kişiye ait her satırda tekrar etmektedir.
\end{itemize}

Caner hoca, özellikle mezunların iş deneyimi örüntüleri ve sektör bazlı
kariyer yollarıyla ilgilenmektedir. Bu bilgiler, bölümün kariyer
danışmanlığı hizmetlerini iyileştirmek için kullanılacaktır. Artık
verileri analiz etmeye hazırsınız!

\section{R ile Mezun Verilerini Analiz
Etme}\label{r-ile-mezun-verilerini-analiz-etme}

\subsection*{Gerekli kütüphaneleri içe
aktar}\label{gerekli-kuxfctuxfcphaneleri-iuxe7e-aktar}
\addcontentsline{toc}{subsection}{Gerekli kütüphaneleri içe aktar}

Aşağıdaki 3 paket sonraki analizler için gereklidir. Ekstra paketler
kullanmak isterseniz buraya ekleyebilirsiniz. Bilgisayarınızda olmayan
paketleri içe aktarmadan önce \texttt{install.packages("...")} ile
yükleyiniz.

\begin{Shaded}
\begin{Highlighting}[]
\FunctionTok{library}\NormalTok{(dplyr)}
\FunctionTok{library}\NormalTok{(ggplot2)}
\FunctionTok{library}\NormalTok{(lubridate)}
\end{Highlighting}
\end{Shaded}

\subsection*{EMU Hacettepe mezun verisini R'a
yükleyin}\label{emu-hacettepe-mezun-verisini-ra-yuxfckleyin}
\addcontentsline{toc}{subsection}{EMU Hacettepe mezun verisini R'a
yükleyin}

Caner hoca veri setini size \texttt{.RData} dosyası
(\texttt{myemu\_data.RData}) olarak vermiştir. Veri dosyasına
\hyperref[0]{www.hadi.hacettepe.edu.tr}'de, ders sayfamızda, sınav
gönderiminden erişebilirsiniz. Dosyayı, R quarto dosyanızla aynı klasöre
koyunuz.

\begin{Shaded}
\begin{Highlighting}[]
\CommentTok{\# Veri setiniz myemu\_data değişkenine data frame olarak yüklenecektir.}
\FunctionTok{load}\NormalTok{(}\StringTok{"myemu\_data.RData"}\NormalTok{)}
\end{Highlighting}
\end{Shaded}

\subsection{Görev 1. Temel Veri
Keşfi}\label{guxf6rev-1.-temel-veri-keux15ffi}

Isınmak için basit görevlerle başlıyorsunuz.

Veri setinde kaç tane benzersiz (unique) şirket (\texttt{work\_company})
var? (sonucu yazdırın)

\begin{Shaded}
\begin{Highlighting}[]
\CommentTok{\# unique\_companies \textless{}{-} l........(u..........(myemu\_data$work\_company))}
\CommentTok{\# print(unique\_companies)}
\end{Highlighting}
\end{Shaded}

Veri setinde kaç tane benzersiz sektör (\texttt{work\_industry}) var?
(sonucu yazdırın)

\begin{Shaded}
\begin{Highlighting}[]
\CommentTok{\# unique\_industries \textless{}{-} l........(u..........(myemu\_data$work\_industry))}
\CommentTok{\# print(unique\_industries)}
\end{Highlighting}
\end{Shaded}

Her bir mezunun kaç farklı iş deneyimi kaydı olduğunu hesaplayınız.
Sonuçları \texttt{jobs\_per\_graduate} isimli bir veri çerçevesinde
(\texttt{data.frame}) saklayınız. Bu veri çerçevesinde
\texttt{student\_id} ve \texttt{job\_count} (iş sayısı) sütunları
bulunmalıdır.

\textbf{İpucu 1:} \texttt{group\_by()} ile \texttt{student\_id}'ye göre
gruplayın.

\textbf{İpucu 2:} \texttt{summarise()} ile \texttt{n()} fonksiyonunu
kullanarak her öğrencinin iş sayısını hesaplayın.

\begin{Shaded}
\begin{Highlighting}[]
\CommentTok{\# jobs\_per\_graduate \textless{}{-} myemu\_data |\textgreater{}}
\CommentTok{\#   g..........(..............) |\textgreater{}}
\CommentTok{\#   s..........(job\_count = ..........)}
\end{Highlighting}
\end{Shaded}

Ortalama iş deneyimi sayısını ve en fazla iş deneyimine sahip mezunun
kaç iş deneyimi olduğunu yazdırın.

\begin{Shaded}
\begin{Highlighting}[]
\CommentTok{\# print(mean(jobs\_per\_graduate$....................))}
\CommentTok{\# print(max(jobs\_per\_graduate$....................))}
\end{Highlighting}
\end{Shaded}

\subsection{Görev 2. En Çok Mezun İstihdam Eden
Şirketler}\label{guxf6rev-2.-en-uxe7ok-mezun-istihdam-eden-ux15firketler}

Caner Hoca, hangi şirketlerin bölüm mezunlarını en çok istihdam ettiğini
merak etmektedir. En çok mezun istihdam eden 10 şirketi tespit edecek ve
bunu görselleştireceksiniz.

\subsubsection{Şirket Başına İstihdam Sayısını
Hesapla}\label{ux15firket-baux15fux131na-istihdam-sayux131sux131nux131-hesapla}

\textbf{Not:} Bir şirket aynı kişiyi farklı pozisyonlarda çalıştırmış
olabilir. Bu durumları ayrı istihdam kayıtları olarak
değerlendirebilirsiniz.

\textbf{İpucu 1:} \texttt{count()} fonksiyonu ile \texttt{work\_company}
sütunundaki kayıt sayılarını hesaplayın.

\textbf{İpucu 2:} \texttt{arrange()} ile azalan sırada düzenleyin.

\textbf{İpucu 3:} \texttt{slice\_head()} ile ilk 10 şirketi seçin.

\begin{Shaded}
\begin{Highlighting}[]
\CommentTok{\# top10\_companies \textless{}{-} myemu\_data |\textgreater{}}
\CommentTok{\#   c........(.............) |\textgreater{}}
\CommentTok{\#   a........(d...........) |\textgreater{}}
\CommentTok{\#   s............(n = ..........)}
\CommentTok{\#}
\CommentTok{\# print(top10\_companies)}
\end{Highlighting}
\end{Shaded}

\subsubsection{En Çok İstihdam Eden 10 Şirket İçin Yatay Çubuk Grafiği
Oluştur}\label{en-uxe7ok-istihdam-eden-10-ux15firket-iuxe7in-yatay-uxe7ubuk-grafiux11fi-oluux15ftur}

Şirket isimlerinin uzun olabileceğini göz önünde bulundurarak, yatay bir
çubuk grafiği oluşturun. Şirketleri istihdam sayısına göre sıralayın.

\textbf{İpucu 1:} \texttt{geom\_bar(stat\ =\ "identity")} kullanın.

\textbf{İpucu 2:} \texttt{reorder()} ile şirketleri \texttt{n} değerine
göre sıralayın.

\textbf{İpucu 3:} Yatay çubuk grafiği için \texttt{y} eksenine şirket
isimlerini, \texttt{x} eksenine sayıları yerleştirin.

\begin{Shaded}
\begin{Highlighting}[]
\CommentTok{\# top10\_companies |\textgreater{}}
\CommentTok{\#   ggplot(aes(x = ..........., y = reorder(............, ...........))) +}
\CommentTok{\#   g..........(stat = "identity", fill = "steelblue") +}
\CommentTok{\#   labs(title = "Top 10 Companies by Number of Graduate Hires",}
\CommentTok{\#        x = "Number of Hires",}
\CommentTok{\#        y = "Company") +}
\CommentTok{\#   theme\_minimal()}
\end{Highlighting}
\end{Shaded}

\subsection{Görev 3. Mezuniyet Yılına Göre CGPA
Analizi}\label{guxf6rev-3.-mezuniyet-yux131lux131na-guxf6re-cgpa-analizi}

Caner Hoca, yıllar içinde mezunların akademik performansının nasıl
değiştiğini merak etmektedir. Mezuniyet yılına göre CGPA dağılımını
inceleyeceksiniz.

\subsubsection{\texorpdfstring{\texttt{graduation\_year} sütunu
ekleyin.}{graduation\_year sütunu ekleyin.}}\label{graduation_year-suxfctunu-ekleyin.}

Veri setinde \texttt{graduation\_date} sütunu \texttt{GG/AA/YYYY}
formatındadır. Ancak size sadece yıl bilgisi gereklidir.

\begin{itemize}
\item
  İlk olarak, \textbf{\texttt{lubridate}} kütüphanesinden
  \textbf{\texttt{dmy()}} fonksiyonunu kullanarak
  \textbf{\texttt{graduation\_date}}'i tarihe dönüştürün.
\item
  Sonra, \textbf{\texttt{year()}} fonksiyonu ile yılı çıkarın.
\item
  Son olarak \textbf{\texttt{graduation\_year}} adında yeni bir sütun
  olarak veri çerçevenize ekleyin.
\item
  Sonraki analizlerde tutarlılık için
  \textbf{\texttt{graduation\_year}}'ı faktör tipi bir değişken (sütun)
  olarak saklayınız.
\end{itemize}

\begin{Shaded}
\begin{Highlighting}[]
\CommentTok{\# myemu\_data \textless{}{-} myemu\_data |\textgreater{}}
\CommentTok{\#   mutate(graduation\_year = as.factor(year(dmy(..........................))))}
\end{Highlighting}
\end{Shaded}

\subsubsection{Yıl Bazında Özet İstatistikleri
Hesapla}\label{yux131l-bazux131nda-uxf6zet-istatistikleri-hesapla}

Her mezuniyet yılı için mezun sayısı, medyan CGPA, minimum CGPA ve
maksimum CGPA hesaplayın.

\textbf{İpucu 1:} Her mezunun sadece ilk görünümünü korumak için
\textbf{\texttt{distinct()}} fonksiyonunu kullanın (\texttt{student\_id}
sütununa göre, \texttt{keep\_all\ =\ TRUE} ile).

\textbf{İpucu 2:} \texttt{group\_by()} ile \texttt{graduation\_year}'a
göre gruplayın.

\textbf{İpucu 3:} \texttt{summarise()} ile gerekli istatistikleri
hesaplayın.

\textbf{İpucu 4:} Mezun sayısını bulmak için \textbf{\texttt{n()}}
fonksiyonunu kullanın.

\begin{Shaded}
\begin{Highlighting}[]
\CommentTok{\# myemu\_data |\textgreater{}}
\CommentTok{\#   d........(.............., .keep\_all = TRUE) |\textgreater{}}
\CommentTok{\#   g........(....................) |\textgreater{}}
\CommentTok{\#   s..........(}
\CommentTok{\#     count\_graduates = ............,}
\CommentTok{\#     median\_cgpa = m...........(....................),}
\CommentTok{\#     min\_cgpa = m....(....................),}
\CommentTok{\#     max\_cgpa = m....(.....................)}
\CommentTok{\#   )}
\end{Highlighting}
\end{Shaded}

\subsubsection{Yıl Bazında Mezun Sayısı Çubuk Grafiği
Oluştur}\label{yux131l-bazux131nda-mezun-sayux131sux131-uxe7ubuk-grafiux11fi-oluux15ftur}

Her yılda kaç mezun olduğunu gösteren bir çubuk grafiği oluşturun.

\textbf{İpucu 1:} Her mezunun yalnızca bir kez temsil edilmesi için
\texttt{distinct()} kullanın.

\textbf{İpucu 2:} \texttt{geom\_bar()} kullanın (sayımı ggplot otomatik
yapacaktır).

\textbf{İpucu 3:} Yılları farklı renklerle göstermek için \texttt{fill}
argümanını kullanın.

\begin{Shaded}
\begin{Highlighting}[]
\CommentTok{\# myemu\_data |\textgreater{}}
\CommentTok{\#   d........(.................., .keep\_all = TRUE) |\textgreater{}}
\CommentTok{\#   ggplot(aes(x = .................., fill = ..................)) +}
\CommentTok{\#   g..........() +}
\CommentTok{\#   labs(title = "Number of Graduates by Year",}
\CommentTok{\#        x = "Graduation Year",}
\CommentTok{\#        y = "Number of Graduates") +}
\CommentTok{\#   theme\_minimal() +}
\CommentTok{\#   theme(axis.text.x = element\_text(angle = 45, hjust = 1))}
\end{Highlighting}
\end{Shaded}

\subsubsection{Her Mezuniyet Yılı İçin Yoğunluk Grafiği (Density Plot)
Çizme}\label{her-mezuniyet-yux131lux131-iuxe7in-youx11funluk-grafiux11fi-density-plot-uxe7izme}

CGPA dağılımını yıllar bazında görebilmek için yoğunluk grafiği
oluşturun. \texttt{facet\_wrap()} ile her yıl için ayrı bir panel
oluşturun.

\textbf{İpucu 1:} Her mezunun yalnızca bir kez temsil edilmesi için
\texttt{distinct()} kullanın.

\textbf{İpucu 2:} \texttt{geom\_density()} kullanın ve
\texttt{fill\ =\ "blue",\ alpha\ =\ 0.5} argümanlarını ekleyin.

\textbf{İpucu 3:} x ekseninde \texttt{graduation\_cgpa} yer alacaktır.

\textbf{İpucu 4:} \texttt{facet\_wrap()} ile \texttt{graduation\_year}'a
göre yüzeyleme yapın.

\begin{Shaded}
\begin{Highlighting}[]
\CommentTok{\# myemu\_data |\textgreater{}}
\CommentTok{\#   d........(................, .keep\_all = TRUE) |\textgreater{}}
\CommentTok{\#   ggplot(aes(x = .....................)) +}
\CommentTok{\#   g...............(fill = "blue", alpha = 0.5) +}
\CommentTok{\#   f............(\textasciitilde{} ..............................) +}
\CommentTok{\#   labs(title = "Density of Graduation CGPA by Year",}
\CommentTok{\#        x = "Graduation CGPA",}
\CommentTok{\#        y = "Density") +}
\CommentTok{\#   theme\_minimal()}
\end{Highlighting}
\end{Shaded}

Yoğunluk grafiklerine dayanarak hangi yılda mezun olanların CGPA
dağılımı en geniş yayılıma sahiptir? Yıllar arasında belirgin bir CGPA
farkı gözlemliyor musunuz?

\begin{longtable}[]{@{}l@{}}
\toprule\noalign{}
Cevap: \\
\midrule\noalign{}
\endhead
\bottomrule\noalign{}
\endlastfoot
Cevabınızı buraya yazınız. \\
\end{longtable}

\subsection{Görev 4: Sektör Bazında CGPA
Karşılaştırması}\label{guxf6rev-4-sektuxf6r-bazux131nda-cgpa-karux15fux131laux15ftux131rmasux131}

Caner Hoca, farklı sektörlerde çalışan mezunların akademik
performanslarının nasıl farklılaştığını merak etmektedir. En çok mezun
istihdam eden sektörlerdeki CGPA dağılımlarını inceleyeceksiniz.

\subsubsection{En Çok İstihdam Kaydına Sahip 8 Sektörü
Belirle}\label{en-uxe7ok-istihdam-kaydux131na-sahip-8-sektuxf6ruxfc-belirle}

\textbf{İpucu 1:} \texttt{count()} ile \texttt{work\_industry}
sütunundaki kayıt sayılarını hesaplayın.

\textbf{İpucu 2:} \texttt{slice\_max()} fonksiyonu ile en yüksek
\texttt{n} değerine sahip 8 sektörü seçin.

\begin{Shaded}
\begin{Highlighting}[]
\CommentTok{\# top8\_industries \textless{}{-} myemu\_data |\textgreater{}}
\CommentTok{\#   c........(.............) |\textgreater{}}
\CommentTok{\#   s............(............., n = ........)}
\CommentTok{\#}
\CommentTok{\# print(top8\_industries)}
\end{Highlighting}
\end{Shaded}

\subsubsection{Bu 8 Sektör İçin CGPA Boxplot Grafiği
Oluştur}\label{bu-8-sektuxf6r-iuxe7in-cgpa-boxplot-grafiux11fi-oluux15ftur}

Sadece yukarıda belirlediğiniz 8 sektöre ait verileri filtreleyerek,
sektör bazında CGPA boxplot grafiği oluşturun. Sektörleri medyan CGPA'ya
göre sıralayın.

\textbf{İpucu 1:} \texttt{filter()} ile
\texttt{work\_industry\ \%in\%\ top8\_industries\$work\_industry}
kullanarak filtreleme yapın.

\textbf{İpucu 2:} \texttt{reorder()} fonksiyonu ile sektörleri medyan
CGPA'ya göre sıralayın.

\begin{Shaded}
\begin{Highlighting}[]
\CommentTok{\# myemu\_data |\textgreater{}}
\CommentTok{\#   f........(work\_industry \%............\% top8\_industries$work\_industry) |\textgreater{}}
\CommentTok{\#   ggplot(aes(x = reorder(..................................), y = ...................., fill = ....................)) +}
\CommentTok{\#   g..............() +}
\CommentTok{\#   labs(title = "CGPA Distribution by Top 8 Industries (Ordered by Median)",}
\CommentTok{\#        x = "Industry",}
\CommentTok{\#        y = "Graduation CGPA") +}
\CommentTok{\#   theme\_minimal() +}
\CommentTok{\#   theme(axis.text.x = element\_text(angle = 45, hjust = 1),}
\CommentTok{\#         legend.position = "none")}
\end{Highlighting}
\end{Shaded}

Sektörler arasında CGPA dağılımı açısından dikkat çekici bir farklılık
var mı? En yüksek medyan CGPA'ya sahip sektör hangisidir?

\begin{longtable}[]{@{}l@{}}
\toprule\noalign{}
Cevap: \\
\midrule\noalign{}
\endhead
\bottomrule\noalign{}
\endlastfoot
Cevabınızı buraya yazınız. \\
\end{longtable}

\subsection{Görev 5: Erasmus Katılımı ve Sektör Tercih
İlişkisi}\label{guxf6rev-5-erasmus-katux131lux131mux131-ve-sektuxf6r-tercih-iliux15fkisi}

Caner Hoca, Erasmus programına katılan ve katılmayan mezunların sektör
tercihlerinde farklılık olup olmadığını merak etmektedir.

\subsubsection{Erasmus Katılım Faktörü
Oluştur}\label{erasmus-katux131lux131m-faktuxf6ruxfc-oluux15ftur}

Veri setinde \texttt{erasmus\_participation} sütunu, Erasmus programına
katılım durumunu göstermektedir. Katılımcılar için bu sütun 1 değerini
alırken, programa katılmayanlar için sütunda NA değeri bulunmaktadır.
\texttt{ifelse()} ve \texttt{is.na()} fonksiyonlarını iç içe kullanarak,
NA değerlerini \texttt{"Hayır"}, NA olmayanları (1'leri) \texttt{"Evet"}
olarak dönüştürün ve sonuçları \texttt{erasmus\_factor} isimli yeni bir
sütunda faktör türünde kaydedin.

\begin{Shaded}
\begin{Highlighting}[]
\CommentTok{\# myemu\_data \textless{}{-} myemu\_data |\textgreater{}}
\CommentTok{\#   m........(erasmus\_factor = factor(i........(i.....(erasmus\_participation), "Hayır", "Evet")))}
\end{Highlighting}
\end{Shaded}

\subsubsection{En Çok İstihdam Kaydına Sahip 6 Sektör İçin Erasmus
Katılım Oranı
Hesapla}\label{en-uxe7ok-istihdam-kaydux131na-sahip-6-sektuxf6r-iuxe7in-erasmus-katux131lux131m-oranux131-hesapla}

Önce en çok istihdam kaydına sahip 6 sektörü belirleyin. Sonra bu
sektörler için Erasmus'a katılan mezunların oranını hesaplayın.

\textbf{İpucu 1:} En çok kayda sahip 6 sektörü belirlemek için önceki
görevdeki yaklaşımı kullanın.

\textbf{İpucu 2:} Bu sektörlerdeki verileri \texttt{filter()} ile
filtreleyin.

\textbf{İpucu 3:} \texttt{group\_by(work\_industry)} ile gruplayın.

\textbf{İpucu 4:} \texttt{summarise()} içinde
\texttt{mean(erasmus\_factor\ ==\ "Evet")} ile oranı hesaplayın.

\begin{Shaded}
\begin{Highlighting}[]
\CommentTok{\# top6\_industries \textless{}{-} myemu\_data |\textgreater{}}
\CommentTok{\#   c........(.................) |\textgreater{}}
\CommentTok{\#   s............(............., n = ........)}
\CommentTok{\#}
\CommentTok{\# erasmus\_by\_industry \textless{}{-} myemu\_data |\textgreater{}}
\CommentTok{\#   f........(work\_industry \%in\% top6\_industries$.................) |\textgreater{}}
\CommentTok{\#   g........(..................) |\textgreater{}}
\CommentTok{\#   s..........(erasmus\_rate = m........(................. == "Evet")) |\textgreater{}}
\CommentTok{\#   a........(d...........(....................))}
\CommentTok{\#}
\CommentTok{\# print(erasmus\_by\_industry)}
\end{Highlighting}
\end{Shaded}

\subsubsection{Sektör Bazında Erasmus Katılım Oranı Çubuk Grafiği
Oluştur}\label{sektuxf6r-bazux131nda-erasmus-katux131lux131m-oranux131-uxe7ubuk-grafiux11fi-oluux15ftur}

Sektörleri Erasmus katılım oranına göre sıralayarak bir çubuk grafiği
oluşturun.

\textbf{İpucu:} \texttt{reorder()} ile sektörleri
\texttt{erasmus\_rate}'e göre sıralayın.

\begin{Shaded}
\begin{Highlighting}[]
\CommentTok{\# erasmus\_by\_industry |\textgreater{}}
\CommentTok{\#   ggplot(aes(x = reorder(..................., ....................), y = ....................)) +}
\CommentTok{\#   g..........(fill = "darkgreen") +}
\CommentTok{\#   labs(title = "Erasmus Participation Rate by Industry (Top 6)",}
\CommentTok{\#        x = "Industry",}
\CommentTok{\#        y = "Erasmus Participation Rate") +}
\CommentTok{\#   theme\_minimal() +}
\CommentTok{\#   theme(axis.text.x = element\_text(angle = 45, hjust = 1)) +}
\CommentTok{\#   scale\_y\_continuous(labels = scales::percent)}
\end{Highlighting}
\end{Shaded}

Erasmus katılım oranları sektörler arasında farklılık gösteriyor mu? Bu
farklılıklara dayanarak ``Erasmus'a katılmak belirli sektörlerde çalışma
olasılığını artırmaktadır'' sonucuna varabilir misiniz? Nedenini
açıklayınız.

\begin{longtable}[]{@{}l@{}}
\toprule\noalign{}
Cevap: \\
\midrule\noalign{}
\endhead
\bottomrule\noalign{}
\endlastfoot
Cevabınızı buraya yazınız. \\
\end{longtable}




\end{document}
